\scanpage{168}
\mspage{85}

\[ \substack{\text{Commutant}\\\text{de }G_0\text{ dans }\unsure{\mathfrak S|\mathscr S}_S}\simeq\operatorname{Norm}_{G_0}(\widetilde U)/\widetilde U \tag*{(35)} \]

donc la donnée d'un sous-groupe $Z$ du commutant revient à la donnée d'un sous-groupe $\widetilde B''$ de $G_0$, \deleted{\unsure{en}} satisfaisant

\[ \widetilde U\subset\widetilde B''\subset\operatorname{Norm}_{G_0}(\widetilde U),\quad\text{avec }Z= \tag*{(36)} \]

\inserted{i.e. $\widetilde B''$ un sous-groupe de $G_0$, contenant $\widetilde U$ comme sous-groupe \uline{distingué} de sorte qu'on a un groupe quotient $Z=\widetilde B''/\widetilde U$}

en termes de (26) on aura

\[ \widetilde B''=\{g\in G_0\mid gs_0\in Z.s_0\}.\tag*{(37)} \]

La condition que $Z$ commute non seulement à l'action de $G_0$, mais aussi à l'action de $G$ \deleted{tout} entier, ou ce qui revient au $=$ de $B'$ (puisque $G=B'.G_0$ \} produit de sous-ens. de $G$) s'explicite aisément par \inserted{l'une} des conditions équivalentes

\[
\begin{aligned}
\text{a)}\quad&\widetilde B''\text{ normalise le sous groupe $B'$ de $G$, ou encore :}\\
\text{b)}\quad&B'\text{ normalise $\widetilde B''$ (pour l'action (33)) et l'action}\\
&\text{dans le quotient $\widetilde B''/\widetilde U$ est triviale.}
\end{aligned}\tag*{(38)}
\]

Ici c'est bien sûr la forme 38 b) qui sera la plus maniable pour notre propos, \deleted{\unsure{en|$Z$}} s'exprimant directement en termes des données (39) et (36). Nous pouvons introduire alors le \inserted{ss-}groupe $B$ de $G$ engendré par $\widetilde B'',B'$, qui s'insère dans le diagramme

\diagram{\begin{tikzcd}[ampersand replacement=\&,column sep=normal,row sep=small]
(39)\quad\&G_0\arrow[r,dash,"Z'"]\arrow[d,dash]\&G\arrow[d,dash]\\
\&\widetilde B''\arrow[r,dash,"Z'"]\arrow[d,dash,"Z"']\&B\arrow[d,dash,"Z"]\\
\&\widetilde U\arrow[r,dash,"Z'"]\&B'
\end{tikzcd}}

où on a marqué au dessus des flèches d'\unsure{inclusions|inclusion}, d'un

% archive sheet number: 167
